\documentclass[11pt,a4paper]{article}
% preamble.tex — estilo común de todos los apuntes Froth.
% Cada apunte hace % preamble.tex — estilo común de todos los apuntes Froth.
% Cada apunte hace % preamble.tex — estilo común de todos los apuntes Froth.
% Cada apunte hace \input{preamble} justo después de \documentclass.
\usepackage[margin=2.4cm]{geometry}
\usepackage{xcolor}
\definecolor{litblue}{RGB}{31,79,121}
\definecolor{litgreen}{RGB}{27,94,32}
\definecolor{litorange}{RGB}{230,81,0}
\usepackage{parskip}
\usepackage{titlesec}
\titleformat{\section}{\Large\bfseries\color{litblue}}{\thesection}{0.6em}{}
\titleformat{\subsection}{\large\bfseries\color{litblue}}{\thesubsection}{0.6em}{}
\usepackage{tcolorbox}
\usepackage{fancyhdr}
\pagestyle{fancy}
\fancyhf{}
\fancyhead[L]{\small\color{litblue}Froth — Apuntes de ML}
\fancyhead[R]{\small\thepage}
\renewcommand{\headrulewidth}{0.4pt}
\usepackage[colorlinks=true,linkcolor=litblue,urlcolor=litblue]{hyperref}

% Cabecera de cada apunte: \cabecera{numero}{titulo}
\newcommand{\cabecera}[2]{%
  \begin{tcolorbox}[colback=litblue,coltext=white,colframe=litblue,arc=2mm]
    {\Large\bfseries Apunte #1 — #2}\\[3pt]
    {\small Proyecto Froth · aprender Machine Learning construyendo}
  \end{tcolorbox}\vspace{0.8em}}

% Cajas temáticas
\newtcolorbox{concepto}[1]{colback=blue!4,colframe=litblue,title={Concepto: #1},arc=1.5mm}
\newtcolorbox{practica}{colback=green!5,colframe=litgreen,title={Pruébalo tú (teclado en mano)},arc=1.5mm}
\newtcolorbox{ojo}{colback=orange!8,colframe=litorange,title={Ojo / error típico},arc=1.5mm}
\newtcolorbox{resumen}{colback=gray!8,colframe=gray!60,title={En una frase},arc=1.5mm}

\newcommand{\codigo}[1]{\texttt{#1}}
 justo después de \documentclass.
\usepackage[margin=2.4cm]{geometry}
\usepackage{xcolor}
\definecolor{litblue}{RGB}{31,79,121}
\definecolor{litgreen}{RGB}{27,94,32}
\definecolor{litorange}{RGB}{230,81,0}
\usepackage{parskip}
\usepackage{titlesec}
\titleformat{\section}{\Large\bfseries\color{litblue}}{\thesection}{0.6em}{}
\titleformat{\subsection}{\large\bfseries\color{litblue}}{\thesubsection}{0.6em}{}
\usepackage{tcolorbox}
\usepackage{fancyhdr}
\pagestyle{fancy}
\fancyhf{}
\fancyhead[L]{\small\color{litblue}Froth — Apuntes de ML}
\fancyhead[R]{\small\thepage}
\renewcommand{\headrulewidth}{0.4pt}
\usepackage[colorlinks=true,linkcolor=litblue,urlcolor=litblue]{hyperref}

% Cabecera de cada apunte: \cabecera{numero}{titulo}
\newcommand{\cabecera}[2]{%
  \begin{tcolorbox}[colback=litblue,coltext=white,colframe=litblue,arc=2mm]
    {\Large\bfseries Apunte #1 — #2}\\[3pt]
    {\small Proyecto Froth · aprender Machine Learning construyendo}
  \end{tcolorbox}\vspace{0.8em}}

% Cajas temáticas
\newtcolorbox{concepto}[1]{colback=blue!4,colframe=litblue,title={Concepto: #1},arc=1.5mm}
\newtcolorbox{practica}{colback=green!5,colframe=litgreen,title={Pruébalo tú (teclado en mano)},arc=1.5mm}
\newtcolorbox{ojo}{colback=orange!8,colframe=litorange,title={Ojo / error típico},arc=1.5mm}
\newtcolorbox{resumen}{colback=gray!8,colframe=gray!60,title={En una frase},arc=1.5mm}

\newcommand{\codigo}[1]{\texttt{#1}}
 justo después de \documentclass.
\usepackage[margin=2.4cm]{geometry}
\usepackage{xcolor}
\definecolor{litblue}{RGB}{31,79,121}
\definecolor{litgreen}{RGB}{27,94,32}
\definecolor{litorange}{RGB}{230,81,0}
\usepackage{parskip}
\usepackage{titlesec}
\titleformat{\section}{\Large\bfseries\color{litblue}}{\thesection}{0.6em}{}
\titleformat{\subsection}{\large\bfseries\color{litblue}}{\thesubsection}{0.6em}{}
\usepackage{tcolorbox}
\usepackage{fancyhdr}
\pagestyle{fancy}
\fancyhf{}
\fancyhead[L]{\small\color{litblue}Froth — Apuntes de ML}
\fancyhead[R]{\small\thepage}
\renewcommand{\headrulewidth}{0.4pt}
\usepackage[colorlinks=true,linkcolor=litblue,urlcolor=litblue]{hyperref}

% Cabecera de cada apunte: \cabecera{numero}{titulo}
\newcommand{\cabecera}[2]{%
  \begin{tcolorbox}[colback=litblue,coltext=white,colframe=litblue,arc=2mm]
    {\Large\bfseries Apunte #1 — #2}\\[3pt]
    {\small Proyecto Froth · aprender Machine Learning construyendo}
  \end{tcolorbox}\vspace{0.8em}}

% Cajas temáticas
\newtcolorbox{concepto}[1]{colback=blue!4,colframe=litblue,title={Concepto: #1},arc=1.5mm}
\newtcolorbox{practica}{colback=green!5,colframe=litgreen,title={Pruébalo tú (teclado en mano)},arc=1.5mm}
\newtcolorbox{ojo}{colback=orange!8,colframe=litorange,title={Ojo / error típico},arc=1.5mm}
\newtcolorbox{resumen}{colback=gray!8,colframe=gray!60,title={En una frase},arc=1.5mm}

\newcommand{\codigo}[1]{\texttt{#1}}

\usepackage{tikz}
\begin{document}
\cabecera{08}{Similitud coseno: medir qué tan parecidos son dos papers}

\section{El problema}

Ya tienes 126 puntos (vectores) en el mapa del significado. Pregunta natural: dados dos
papers, ¿cómo pone la computadora un número a ``qué tan parecidos son''? Con la
\textbf{similitud coseno}.

\section{La idea: el ángulo entre dos flechas}

Cada vector es una \textbf{flecha} que sale del origen y apunta hacia la posición del paper.
Para comparar dos papers miramos el \textbf{ángulo} entre sus flechas:

\begin{center}
\begin{tikzpicture}[scale=1.0]
  % caso parecido
  \draw[->,gray] (0,0) -- (3.2,0);
  \draw[->,gray] (0,0) -- (0,3.2);
  \draw[->,litblue,very thick] (0,0) -- (2.6,2.2);
  \draw[->,litgreen,very thick] (0,0) -- (2.9,1.6);
  \draw (1.0,0.62) arc (32:40:1.2);
  \node at (1.7,0.55) {\small ángulo chico};
  \node[align=center] at (1.6,-0.8) {\small\textbf{parecidos}\\[-2pt]\small coseno $\approx 1$};
  % separacion
  \node at (4.6,1.5) {\Large $\Rightarrow$};
  % caso distinto
  \begin{scope}[xshift=6cm]
    \draw[->,gray] (0,0) -- (3.2,0);
    \draw[->,gray] (0,0) -- (0,3.2);
    \draw[->,litblue,very thick] (0,0) -- (2.9,0.5);
    \draw[->,litorange,very thick] (0,0) -- (0.5,2.9);
    \node[align=center] at (1.6,-0.8) {\small\textbf{distintos}\\[-2pt]\small coseno $\approx 0$};
  \end{scope}
\end{tikzpicture}
\end{center}

\begin{concepto}{similitud coseno}
Es el \textbf{coseno del ángulo} entre dos vectores. Da un número fácil de leer:
\begin{itemize}
  \item Flechas casi alineadas (ángulo pequeño) $\rightarrow$ coseno \textbf{cercano a 1}
        $\rightarrow$ mismo tema.
  \item Flechas perpendiculares (90$^\circ$) $\rightarrow$ coseno \textbf{cercano a 0}
        $\rightarrow$ sin relación.
\end{itemize}
Para embeddings de texto el valor suele caer entre 0 y 1.
\end{concepto}

\begin{ojo}
¿Por qué el \emph{ángulo} y no la \emph{distancia} entre las puntas? Porque lo que importa
es \textbf{hacia dónde apunta} el significado, no qué tan larga es la flecha. Dos papers del
mismo tema, uno con abstract largo y otro corto, tendrán flechas de distinto ``tamaño'' pero
apuntando al mismo lado: el ángulo los reconoce como iguales; la distancia recta se dejaría
engañar por el tamaño.
\end{ojo}

\section{La prueba con tus papers (números reales)}

Tomamos el paper \emph{``Enrichment of Lepidolite via Flotation''} y calculamos su coseno
contra los otros 125. Esto salió:

\begin{center}
\begin{tabular}{lll}
\textbf{Coseno} & \textbf{Paper} & \textbf{¿Tiene sentido?} \\
\hline
0.942 & Selective Inhibition ... Flotation of Lepidolite & sí: mismo tema \\
0.941 & Froth Flotation of Lepidolite --- A Review & sí: mismo tema \\
0.932 & Flotation Collectors for Lepidolite Mineral & sí: mismo tema \\
\hline
0.511 & ...Radio Resources for Multicell Mobile & no: telecomunicaciones \\
0.507 & Adrienne Rich and the Poetry... & no: poesía \\
0.495 & Theodore Beza... Peace in France 1572 & no: historia \\
\end{tabular}
\end{center}

Los papers de flotación de lepidolita salieron con coseno $\sim$0.94 (casi idénticos en
tema), y el ruido off-topic que se coló en la Fase 1 quedó en $\sim$0.50 (lejos). El sistema
los ordenó bien \textbf{sin saber nada de minería}: solo por el significado de los abstracts.

\begin{concepto}{la herramienta: \codigo{cosine\_similarity} de scikit-learn}
No calculamos el ángulo a mano. \codigo{sklearn.metrics.pairwise.cosine\_similarity} toma
vectores y devuelve los cosenos. Le pasamos un paper contra la matriz de los 126 y obtuvimos
las 126 similitudes de un golpe.
\end{concepto}

\section{Para qué sirve esto}

Este número es el ladrillo del \textbf{buscador semántico} (paso siguiente): para encontrar
los papers más parecidos a una frase, la convertimos en vector y nos quedamos con los de
\emph{mayor} coseno. Y en la Fase 3, la misma idea de ``quién está cerca de quién'' es la que
forma las burbujas de subtemas.

\begin{resumen}
La similitud coseno mide el ángulo entre dos vectores: $\approx 1$ mismo tema, $\approx 0$ sin
relación. Con tus papers, flotación-vs-flotación dio 0.94 y flotación-vs-ruido 0.50: funciona.
\end{resumen}

\end{document}
